\documentclass[11pt,a4paper]{article}

\usepackage{../riddle}
\usepackage[maxalphanames=99, maxnames=99, backend=bibtex, style=alphabetic, sorting=ynt]{biblatex}
\addbibresource{stage3A.bib}
\DeclareSymbolFont{yhlargesymbols}{OMX}{yhex}{m}{n} 
\DeclareMathAccent{\yhwidehat}{\mathord}{yhlargesymbols}{"62}
\renewcommand{\headrulewidth}{1pt} 
\renewcommand{\footrulewidth}{1pt}	
\fancyhead[C]{}
\fancyhead[L]{}
\fancyhead[R]{}
\fancyfoot[C]{\thepage} 
\fancyfoot[L]{Sacha Ben-Arous}
\fancyfoot[R]{E.N.S Paris-Saclay}
\usepackage{tikz-cd}

% Majda-Bertozzi avec fiches -> Eggers-Fontelos -> Tao chap1 et annexe -> papier et articles connexes 

% https://www.youtube.com/watch?v=mw_hhOqKDx4 (T.B. intro papier 2D)
% https://vimeo.com/539376697/bbc4763e80 (T.B. intro papier 3D)
% https://www.youtube.com/watch?app=desktop&v=JRBwdVFmcE4 (T.B. exemple conf avec op linéarisé) (en chercher d'autres)


%TODO SE REFERER AU DOCUMENT DE F PI !!!!!!!!!


\begin{document}

\begin{titlepage}
\centering
\vspace*{0.5cm}
{\large École Normale Supérieure Paris-Saclay \par}
{\large New York University, Courant Institute \par}
\vspace{1cm}
{\large Research internship conducted as part of the Master 1 Hadamard \par}
{\large Supervised by Tristan Buckmaster\par}

\vspace{3cm}
{\Huge \textbf{Self-similar singularity formation in fluids}\par}
\vspace{1cm}

{\Large {Sacha Ben-Arous}\par}
    \vfill 
     \begin{tabular}{p{0.45\linewidth} p{0.45\linewidth}}
        \raisebox{-1.6cm}{\includegraphics[width=7cm]{logo_accueil.png}} & 
        \hfill \includegraphics[scale=0.25]{logo_origine.jpg}
    \end{tabular}

\end{titlepage}
{\centering
\section*{Compte-rendu de stage}
}
D'Avril à Août 2025, j'ai effectué mon stage de Master 1 à l'Université de New York, sous la direction de Tristan Buckmaster. Ce stage avait pour objectif à la fois de s'initier à un domaine de recherche (très) actif, et de découvrir le fonctionnement du monde de la recherche. Le compte-rendu qui suit décrit mon expérience globale au cours de ce stage. \\

Tout d'abord, ce qui m'a attiré vers ce stage est la possibilité de d'étudier un problème central des mathématiques : le caractère bien-posé de certaines équations de la mécanique des fluides, sous la direction d'un chercheur ayant apporté des contributions significatives à ce domaine : Tristan Buckmaster a participé à la résolution de la conjecture d'Onsager, avec P. Isett et V. Vicol. De plus, le Courant Institute où j'ai effectué mon stage à une grande tradition de l'étude théorique et numérique de la mécanique des fluides, grâce à des mathématiciens illustres tels que P. Lax, L. Nirenberg, A. Majda, C. Morawetz, et plus récemment avec T. Buckmaster, V. Vicol, J. Shatah. \\

Par ailleurs, j'ai dès le début de mon stage été marqué par la pluridisciplinarité présente dans ce laboratoire, grâce au large spectre de disciplines représentées. Un nombre colossal d'exposés avait lieu durant le premier mois de stage (au moins un par jour), dont les sujets variaient de l'étude numérique appliquée de la dynamique des océans jusqu'à l'utilisation de matrices aléatoires pour étudier la fonction zêta de Riemann, en passant par de la théorie de la mesure géométrique, avec la fameuse résolution de la conjecture de Kakeya par Hong Wang en Mars dernier. \\
Un autre fait marquant a été d'assister à la conférence célébrant la retraite de S.R.S. Varadhan, monument des probabilités modernes, ayant passé 60 ans (!!!) au Courant Institute, en ayant été deux fois son directeur. Ce dernier est à l'origine de l'autre tradition majeure de Courant : l'étude et le développement de la théorie des probabilités modernes, en particulier des phénomènes de grandes déviations, ainsi que de l'étude des processus de diffusion. Cette conférence a été l'occasion pour moi de découvrir des facettes souvent négligées dans les description des grandes carrières scientifiques : Varadhan a bien sur été extrêment prolifique sur le plan des découvertes mathématiques, mais il s'est aussi et surtout massivement investi dans la construction d'une culture de la recherche en probabilité qui porte encore ses fruits aujourd'hui, par exemple avec l'étude des matrices aléatoires citée ci-dessus, mais aussi avec l'intelligence artificielle, qui est aujourd'hui sans conteste le domaine le plus actif de la recherche mathématique, et dans lequel le Courant Institute est un lieu d'avant garde. \\
Finalement, j'ai aussi eu la chance de participer  à la fin de mon stage à une summer school organisée par la National Science Foundation, portant précisemment sur le thème de mon stage, qui se déroulait au département de mathématiques de Princeton. Cela m'a permit d'avoir d'être exposé à des travaux actuels, d'avoir de très nombreuses discussions avec les différents participants, qu'ils soient professeurs ou élèves, et d'observer comment se coordonne un groupe de recherche en établissant une stratégie à long terme pour résoudre un problème important. \\

Concernant le contenu de mon stage, j'ai été initié à une stratégie très générale employée par Tristan Buckmaster et ses collaborateurs afin d'étudier l'existence de solutions régulières aux équations de la mécanique des fluides, telles que l'équation d'Euler ou l'équation de Navier-Stokes (incompressibles). Ce problème est bien sur très célèbre, car il est mis à prix par le Clay Institute pour un million de dollars, David Hilbert l'ayant décrit comme l'un des problèmes devant guider les recherches des mathématiciens au cours de notre millénaire. Ce problème peut à première vu sembler dénué d'interêt aujourd'hui, tant les modélisations utilisées dans des domaines appliquées tels que la météorologie ou l'étude de la dynamique des océans sont infiniments plus compliqués. Cependant, ce qui motive une telle question n'est pas d'évaluer la pertinence des équations d'Euler ou de Navier-Stokes pour décrire la réalité, mais est plutôt de permettre, via l'abstraction induite par le cadre mathématique, de mettre en évidence les potentielles incohérences et défauts de nos modélisations intuives des fluides conduisant à des absurdités physiques, et donc de repenser à un niveau fondamental notre conception et appréhension des fluides. Le premier résultat majeur sur ce sujet a été publié par Jean Leray en 1934 \cite{leray1934} dans lequel est établit l'existence globale de solutions faibles, "turbulentes", aux équations de Navier-Stokes, aujourd'hui appelées solutions de Leray-Hopf. Ce sujet a été largement étudié depuis, et une avancée importante récente est l'existence de données initiales lisses qui conduisent à une singularité pour l'équation d'Euler à frontière cylindrique, publiée par Chen et Hou en 2023 dans une série d'articles \cite{chen_hou1},\cite{chen_hou2}, et qui utilise une stratégie semblable à celle décrite dans ce rapport. \\
Il faut aussi noter le récent progrès dans la résolution du 6ème problème de Hilbert, qui concerne la dérivation probabiliste des équations de la physique (en particulier de la mécanique des fluides), obtenu par Deng, Hani, Ma \cite{deng_hani_ma_2025_hilbert6}. Ce résultat s'inscrit dans une série de travaux menés par Bodineau, Gallagher, Golse, Saint-Raymond, et offre une approche potentiellement différente de la vision analytique classique, car elle s'appuie sur des outils probabiliste, et est motivée physiquement par la description atomique, moléculaire, des fluides. Cette vision probabiliste des fluides comme étant la limite d'un ensemble de particules en intéraction ne sera cependant pas abordée ici. \\

D'un point de vue pratique, j'ai d'abord commencé par apprendre les résultats classiques de la mécanique des fluides et de l'auto-similarité dans des livres tels que \cite{majda2001vorticity}, \cite{eggers2015singularities}, \cite{tao2006dispersive}, j'ai ensuite lu les articles \cite{buckmaster2019formation},\cite{buckmaster2022imploding} afin de découvrir les différentes technique utilisées pour résoudre le type de problème auquel j'étais confronté. J'ai ensuite expérimenté avec les approximations numériques et l'utilisation des bibliothèques Python de calcul formel et de calculs rigoureux, puis j'ai appliqué la méthode générale précedemment mentionnée dans le cas de l'équation de Burgers.

\newpage
{\centering
\section*{Internship summary}
%TODO 
}

\newpage
{\centering
\section*{Introduction}

}
\section{Introduction}

%TODO : Classification AMS
% Le résultat qu'on veut montrer est beaucoup plus fort que simplement cette equation a une donnée initiale bizarre qui engendre une singularité. On prouve qu'il y a un ouvert de données initiales qui génèrent TOUTES la MÊME singularité vue asymptotiquement via le changement de variable auto-similaire ! Ce n'est donc pas un phénomène illusoire, négligeable (cf triadique de cantor pour faire péter Euler), et de plus Tristan semble penser que le comportement pour Burgers est universel, quitte a bien changer les variables. % Deux trucs dans le papier 1 : géométrie de la singularité stable (au sens C^4 pour la géométrie asymptotique auto similaire du choc) ; et purement azymuthale (en se ramenant à la sol 1D) ; et viscosite non nulle et bornee.
% Citer Chen/Hou blow up cylindrical boundary for smooth inital data.
% Interet de self-smilarity : universallité, unifie. Et deja présent chez Leray ?
% On veut du stable car on fait du numerique et donc si instable alors les erreurs numeriques nous emmène loin. (En fait ps sur, gros trouver une vrai def)

% Parler de la vision probabiliste système de particules en intéraction, 6ème problème de Hilbert et article récent (+ filiation Bardos, Bodineau, Gallagher, Golse, Saint-Raymond)

\section{Analysis of shock and self-similar singularities in a fluid}
% Pour le contenu :
% Se servir des présentations de Tristan !
% On prend une EDP, on veut montrer que des solutions forment des singularités en temps fini, stables, non triviales (densité non nulle). 
% On s'intéresse aux singularités self similaire. Expliquer pourquoi et faire la présentation
% Pour cela on a en fait besoin de montrer que quasimment tout reste bornée, c'est contre intuitif, expliquer.
% Faire l'intro DES papiers des doc.
%Faire un résumé des grandes lignes des techniques ci-dessus avant de les mettre en pratique. Inégalité d'énergie pour opérateur linéarisé au début, bootstrap pour partie non linéaire à la fin
% Étape 0 : partir de Burgers et expliquer self-similar.
% Étape 1 : premier papier facile car même pas besoin d'étudier l'op linéarisé car lemme appendice equation de forcing.  Second papier plus dur car 3d et pas Riemann. Cette fois il faut l'inégalité d'énergie mais on peut s'en sortir théoriquement.
% Étape 2 : mon stage. Cette fois il faut mettre en place des méthodes numériques. Quand même facile car on connait très bien l'équation (on a carrément la solution explicite).
% Étape 3 : Maintenant on ne connait même pas la solution autour de laquelle linéariser. Utilisation des PINNs et approximation par gaussienne(RKHS)/Hermite.
% Étape 4 : premier papier. Décrire méthode barrière.
%TODO : Préciser si on est dans le cas incompressible, visqueux, etc ...
%TODO : donner les références des livres
In this section, I will present the concepts of shock formation and self-similar blow-up for a nonlinear partial differential equation, illustrated by the fundamental example of Burgers' equation.


\section{Energy inequality for the linearized operator}
b
%TODO : Où intervient la partie symétrique de l'opérateur ???? (Normalement on peut IPP donc tout est déjà symétrique)

% DIRE QUE LES MODES QU'IL RESTE À LA FIN SONT TRAITÉS DANS LA PARTIE NON LINÉAIRE, EN CHOISISSANT BIEN LA DONNÉE INITIALE.

\subsection{ Quality of the approximation}
%TODO : Faire un première partie sur les calculs de eta et gamma et dire que gamma_i(j)=gamma_j(i)=gamma_0(i-j)=gamma_0(j-i) car paire.
The compact operators that we are studying are of the form: \[K(f)(x) = \int_\R K(x,y)f(y)\mathrm{d}y \]
Assuming that we can bound the growth of the kernel at infinity, a natural choice of a finite rank approximation on a bounded domain $[-A,A]$ would be to use a Riemann sum : \[K_n(f)(x) := \sum_{i=0}^n \delta K(x,y_i)f(y_i)\]
where $\delta := \frac{2A}{n}$ is the integration step and $\begin{cases} y_0=-A \\ y_{i+1}=y_i + \delta \end{cases}$ are the sample points. \\

Now, we want to get a precise bound on the quality of this approximation in order to compute a relevant upper bound in the energy estimate \ref{energy}. \\
We want to use the following result to get a precise bound on the convergence in the operator norm.
\begin{lemma}[Schur test]\label{schur}
Let $K: \R^2 \mapsto \R$ be a square integrable kernel, and $T$ be the operator defined by 
\[ T: L^2(\R) \longrightarrow L^2(\R), \qquad \big(Tf\big)(x) := \int_{\R}K(x,y) f(y)\mathrm{d}y. \] 
Then \[\displaystyle \|T\|_{L^2\mapsto L^2} \leq \sup_{x\in \R} \int_\R |K(x,y)|\mathrm{d}y\times  \sup_{y\in \R} \int_\R |K(x,y)|\mathrm{d}x.\] 
\end{lemma}

Let's now rewrite the operators to apply this lemma. Recall that the Fourier multiplier $(1-\Delta)^\frac{3}{2}$ defines an isometric isomorphism from $H^3(\R)$ to $L^2(\R)$, and also that the Dirac function that evaluates a function at a given point $x$ has a representation in $H^3(\R)$, that we denote $\eta_x$. Now, for a function $f$ in $H^3(\R)$, defining $g:=(1-\Delta)^\frac{3}{2}f$, we have: 
\[K(f)(x) = \int_\R K(x,y)f(y)\mathrm{d}y = \int_\R(1-\Delta_y)^\frac{-3}{2} K(x,y)g(y)\mathrm{d}y\]
and 
\begin{eqnarray*}
K_n(f)(x) &=& \sum_{i=0}^n \delta K(x,y_i)f(y_i) = \sum_{i=0}^n \delta K(x,y_i)\langle f, \eta_{y_i} \rangle_{H^3} =   \sum_{i=0}^n \delta K(x,y_i)\langle g, (1-\Delta_y)^\frac{3}{2}\eta_{y_i} \rangle_{L^2} \\
&=&\int_\R \sum_{i=0}^n \delta K(x,y_i)(1-\Delta_y)^\frac{3}{2}\eta_{y_i}(y)g(y)\mathrm{d}y 
\end{eqnarray*} 

\subsection{Conclusion for the linear part}


% Avec les notes de Gonzalo.
%Prendre f2(0)=0 dans notre cas
%Prendre f dans l'orthogonal de l'image de M_n
%(On peut meme faire mieux en gardant que les espaces propres de vp trop grandes)
%(En fait la bonne façon de faire est de montrer numériquement que l’intégrale décroit. Je peux court-circuiter avec la preuve bootstrap des slides ?)



\section{Bootstrap method for the non-linear part}
c
% La méthode de bootstrap /continuité revient à montrer qu'une zone "barrière" est interdite, et comme le système évolue continuement, il reste forcément d'un des côtés de la barrière, qui dépend juste de l'état initial.
% expliquer pourquoi le damping damp : lutte contre la dérivéée temporelle, et lemme A.1 donne une estimée précise.
% Deux trucs dans le papier : géométrie de la singularité stable (au sens C^4 pour la géométrie asymptotique auto similaire du choc) ; et purement azymuthale (en se ramenant à la sol 1D) ; et viscosite non nulle et bornee.
\newpage
\printbibliography[heading=bibintoc]
\newpage
\appendix
\section{Computation of the kernel}

\subsection{Definitions}

Let $W(x) := \left( -\frac{x}{2} + \left( \frac{1}{27} + \frac{x^2}{4} \right)^{1/2} \right)^{1/3} -\  \left( \frac{x}{2} + \left( \frac{1}{27} + \frac{x^2}{4} \right)^{1/2} \right)^{1/3}$ \\

%TODO be the solution of the self-similar Burgers' equation : ????

and define $Lz := -az - b\frac{\partial z}{\partial x} $, where $a(x):= 1+\frac{W(x)}{x}+\frac{\partial W}{\partial x} (x)$ and $b(x):= \frac{3x}{2} + W(x)$.

\subsection{Computation in Sobolev spaces}
%TODO : Problème de nb de dérivées dans le calcul de la fq ? Prendre la partir sym ? 
%TODO : PARTIE SYMÉTRIQUE 
%TODO :Vraiment faire la forme quadratique

Let $\langle f, g \rangle := \langle f, g \rangle_{L^2} = \displaystyle \int_\R  f(x)g(x)\,dx$ be the usual inner product of $L^2(\R)$.

Let $w,z \in H^k(\R)$ for $k$ large enough. For simplicity, we will denote $\frac{\partial z}{\partial x} := z'$.

\subsubsection{Symetric part in $L^2$ space}
In $L^2$, the symetric part is computed as follows :
\begin{align*}
  \langle Lz, w \rangle_{L^2}  &=   \langle -az - bz', w \rangle =  \langle z, -aw \rangle  +   \langle z, b'w+bw' \rangle    \\
                  &= \langle z, (-a+b')w + bw' \rangle = \langle z, L^*w \rangle
\end{align*}

Thus, $\frac{1}{2}(L+L^*)z = \frac{1}{2}(-az-bz' -az +b'z+bz') \displaystyle =-az+\frac{b'}{2}z$ in $L^2$.

\subsubsection{Quadratic form in $H^1$ space}
In $H^1$, the quadratic form is computed as follows :
\begin{align*}
  \langle Lz, z \rangle_{H^1}   &=   \langle -az - bz', z \rangle +  \langle -a'z-az' - b'z'-bz'', z' \rangle \\
  &=  \langle -az, z \rangle + \langle -bz - a'z, z' \rangle + \langle - az'-b'z', z' \rangle + \langle - bz', z'' \rangle \\
  &=  \langle -az, z \rangle + \langle \frac{1}{2}(b'+a'')z, z \rangle + \langle (-a-b')z', z' \rangle + \langle \frac{1}{2} b' z', z' \rangle \\
  &=  \langle (-a+ \frac{b'}{2}+\frac{a''}{2})z, z \rangle +  \langle (-a-\frac{b'}{2})z', z' \rangle 
\end{align*}

\begin{rmq}
The operator $(Lz)'$ is not defined on $H^1$ as it involves second derivatives of $z$, but it is a classical fact that the quadratic form of an operator as a larger domain than the operator itself.
\end{rmq}

\subsubsection{Quadratic form in $H^2$ space}
In $H^2$, the quadratic form is computed as follows :
\begin{align*}
  \langle (Lz)'', z'' \rangle   &=    \langle -a''z-a'z'-a'z'-az''- b''z'-b'z''-b'z''-bz^{(3)}, z'' \rangle \\
  &= \langle -a''z, z'' \rangle +  \langle (-2a'- b'')z', z'' \rangle +  \langle (-a-2b')z'', z'' \rangle + \langle -bz^{(3)}, z'' \rangle \\
  &= \langle a^{(3)}z+ a''z', z' \rangle +  \langle \frac{1}{2} (2a''+ b^{(3)})z', z' \rangle +  \langle (-a-2b')z'', z'' \rangle + \langle \frac{1}{2} b'z'', z'' \rangle \\
  &= \langle - \frac{1}{2}a^{(4)}z, z \rangle +  \langle 2a''+ \frac{1}{2} b^{(3)})z', z' \rangle +  \langle (-a- \frac{3}{2} b')z'', z'' \rangle 
\end{align*}

Thus, we have in $H^2$ :
\begin{align*}
  \langle Lz, z \rangle_{H^2}   &=  \langle (-a+ \frac{b'}{2}+\frac{a''}{2}-\frac{a^{(4)}}{2})z, z \rangle +  \langle (-a-\frac{b'}{2} +2a''+ \frac{ b^{(3)}}{2})z', z' \rangle +  \langle (-a- \frac{3}{2} b')z'', z'' \rangle 
\end{align*}



\subsubsection{Quadratic form in $H^3$ space}
In $H^3$, the quadratic form is computed as follows :
\begin{align*}
\langle (Lz)^{(3)} , z^{(3)}  \rangle
&= \langle -a'''z - 3a''z' - 3a'z'' - a z^{(3)} - b'''z' - 3b''z'' - 3b'z^{(3)} - b z^{(4)}, z^{(3)} \rangle \\
&= \langle -a'''z, z^{(3)} \rangle + \langle (-3a'' - b''')z', z^{(3)} \rangle + \langle (-3a' - 3b'')z'', z^{(3)} \rangle + \langle (-a - 3b')z^{(3)}, z^{(3)} \rangle \\
&\quad + \langle -b z^{(4)}, z^{(3)} \rangle \\
&= \langle a^{(4)}z + a'''z', z'' \rangle + \langle (3a''' + b^{(4)})z' + (3a'' + b''')z'', z'' \rangle + \langle \tfrac{3}{2}(a'' + b''')z'', z'' \rangle  \\
&\quad + \langle (-a - 3b')z^{(3)}, z^{(3)} \rangle  + \langle \tfrac{1}{2}b' z^{(3)}, z^{(3)} \rangle \\
&= \langle -a^{(5)}z-a^{(4)}z', z' \rangle + \langle  -\frac{1}{2} a^{(4)}z', z' \rangle  + \langle \frac{1}{2}(-3a^{(4)}  - b^{(5)})z', z' \rangle   + \langle (3a'' + b''')z'', z'' \rangle  \\
&\quad + \langle \tfrac{3}{2}(a'' + b''')z'', z'' \rangle + \langle (-a - 3b')z^{(3)}, z^{(3)} \rangle  + \langle \tfrac{1}{2}b' z^{(3)}, z^{(3)} \rangle \\
&= \langle \frac{a^{(6)}}{2}z , z \rangle +  \langle (-3a^{(4)}  - \frac{1}{2} b^{(5)} )z', z' \rangle  +  \langle (\tfrac{9}{2}a'' +\tfrac{5}{2} b^{(3)})z'', z'' \rangle  + \langle (-a - \frac{5}{2} b')z^{(3)}, z^{(3)} \rangle  
\end{align*}

Thus, we have in $H^3$ :
\begin{align*}
  \langle Lz, z \rangle_{H^3}   &=  \langle (-a+ \frac{b'}{2}+\frac{a''}{2}-\frac{a^{(4)}}{2} + \frac{a^{(6)}}{2})z, z \rangle +  \langle (-a-\frac{b'}{2} +2a''+ \frac{ b^{(3)}}{2}-3a^{(4)} - \frac{1}{2} b^{(5)} ))z', z' \rangle \\
  &\quad +  \langle (-a- \frac{3}{2} b'+\tfrac{9}{2}a'' +\tfrac{5}{2} b^{(3)})z'', z'' \rangle + \langle (-a - \frac{5}{2} b')z^{(3)}, z^{(3)} \rangle 
\end{align*}

\subsection{Compact part of the quadratic form}
%TODO Montrer que les fonctions en jeu sont intégrables, L^2, dérivable et toutes les props nécessaires à l'étude ci-dessous.
%TODO Faire les vraies preuves sur le signe (i.e. plus petit que -1/4).
We proved in the previous section that the quadratic form associated with $L$ in $H^3$ is of the form :
\[ \langle Lz, z \rangle_{H^3} =  \langle \varphi_0 z, z \rangle +\langle \varphi_1 z', z' \rangle +\langle \varphi_2 z'', z'' \rangle +\langle \varphi_3 z^{(3)}, z^{(3)} \rangle\]

In the next section, we will show that $\varphi_3$ has a sign and is bounded. This leaves to study the lower order terms, and we will prove that there exists a compact operator $M$ such that 
\[\langle Mz, z \rangle_{H^3} =  \langle \varphi_0 z, z \rangle +\langle \varphi_1 z', z' \rangle +\langle \varphi_2 z'', z'' \rangle\]

Combining those results yield the following energy estimate :
\begin{equation}\label{energy}
\langle Lz, z \rangle_{H^3} \leq -\delta \|z\|_{H^3}^2  + \langle Mz, z \rangle_{H^3} 
\end{equation}
We will use the Fourier transform, with the following convention :
\[
\hat{f}(\xi) := \mathcal{F}(f)(\xi) = \int_{-\infty}^{\infty} f(x) \, e^{-2\pi i x \xi} \, dx
\]
and we will denote \[\mathcal{F}^{-1}(f)(x) := \int_{-\infty}^{\infty} f(\xi) \, e^{2\pi i x \xi} \, d\xi\]
the inverse Fourier transform.

\subsubsection{Base case}
%TODO check si il manque pas une conjuguaison dans ma def du prod L^2 pour appliquer parseval, et vérifier que les calculs ont du sens car z est dérivé une fois de trop à la fin des IPP.
We want to find $M_0$ such that
\begin{equation}\label{c0}
 \langle M_0z, z \rangle_{H^3} =  \langle \varphi_0 z, z \rangle\
\end{equation} 

The Parseval identity gives : \[\int \hat{z}(\xi) \widehat{M_0z}(\xi)(1+\xi^2)^3 \mathrm{d}\xi= \int \hat{z}(\xi)\widehat{\varphi_0z}(\xi)\mathrm{d}\xi \]

Thus, choosing $M_0$ such that $\widehat{M_0z}(\xi)= \frac{1}{(1+\xi^2)^3}\widehat{\varphi_0z}(\xi)$ would give the equality. \\

Defining $\lambda_0(\xi) := \frac{1}{(1+\xi^2)^3}$, this condition is equivalent to : \[\widehat{M_0z}= \widehat{\mathcal{F}^{-1}(\lambda_0)}\widehat{\varphi_0z}=\yhwidehat{\mathcal{F}^{-1}(\lambda_0)*\varphi_0z}\]
i.e. $M_0z = \mathcal{F}^{-1}(\lambda_0)*\varphi_0z$ satisfies \fcref{c0}.

\subsubsection{First order case}

We want to find $M_1$ such that
\begin{equation}\label{c1}
 \langle M_1z, z \rangle_{H^3} =  \langle \varphi_1 z', z' \rangle\
\end{equation} 

Integrating by parts and applying the Parseval identity, we have the equivalence 
\begin{eqnarray*}
\langle M_1z, z \rangle_{H^3} &=& -  \langle \varphi_1' z'+\varphi_1 z'', z \rangle = -  \langle \varphi_1'z, z' \rangle - \langle \varphi_1z, z'' \rangle \\
\Leftrightarrow \int \hat{z}(\xi) \widehat{M_1z}(\xi)(1+\xi^2)^3 \mathrm{d}\xi &=& - \int (2\pi i \xi) \hat{z}(\xi)\yhwidehat{\varphi_1'z}(\xi)\mathrm{d}\xi + \int (4\pi^2 \xi^2)\hat{z}(\xi)\yhwidehat{\varphi_1z}(\xi)\mathrm{d}\xi \\
\Leftrightarrow \int \hat{z}(\xi) \widehat{M_1z}(\xi)(1+\xi^2)^3 \mathrm{d}\xi &=&  \int\hat{z} \left[-(2\pi i \xi) \yhwidehat{\varphi_1'z}(\xi) + (4\pi^2 \xi^2)\hat{z}(\xi)\yhwidehat{\varphi_1z}(\xi) \right]\mathrm{d}\xi 
\end{eqnarray*}
Defining $\lambda_1(\xi):= - \frac{2\pi i \xi}{(1+\xi^2)^3 }$ and $\lambda_2(\xi) := \frac{4\pi^2 \xi^2}{(1+\xi^2)^3 }$, we have that 
\[M_1z :=  \left(\mathcal{F}^{-1}(\lambda_1)*\varphi_1'z\right)+\left(\mathcal{F}^{-1}(\lambda_2)*\varphi_1z\right) \]
satisfies \fcref{c1}.


\subsubsection{Second order case}

We want to find $M_2$ such that
\begin{equation}\label{c2}
 \langle M_2z, z \rangle_{H^3} =  \langle \varphi_2 z'', z'' \rangle\
\end{equation} 

Integrating by parts twice and applying the Parseval identity, we have the equivalence 
\begin{eqnarray*}
\langle M_2z, z \rangle_{H^3} &=&  \langle \varphi_2''z'' +2\varphi_2'z^{(3)}+\varphi_2z^{(4)}, z \rangle = \langle \varphi_2''z, z'' \rangle + \langle 2 \varphi_2'z, z^{(3)} \rangle + \langle \varphi_2z, z^{(4)} \rangle \\
\Leftrightarrow \int \hat{z}(\xi) \widehat{M_2z}(\xi)(1+\xi^2)^3 \mathrm{d}\xi &=& -  \int (4\pi^2 \xi^2)\hat{z}(\xi)\yhwidehat{\varphi_2''z}(\xi)\mathrm{d}\xi - \int (i16\pi^3\xi^3)\hat{z}(\xi)\yhwidehat{\varphi_2'z}(\xi)\mathrm{d}\xi + \int (16\pi^4\xi^4)\hat{z}(\xi)\yhwidehat{\varphi_2z}(\xi)\mathrm{d}\xi \\
\Leftrightarrow \int \hat{z}(\xi) \widehat{M_2z}(\xi)(1+\xi^2)^3 \mathrm{d}\xi &=&  \int\hat{z} \left[-(4\pi^2 \xi^2) \yhwidehat{\varphi_2''z}(\xi) - (i16\pi^3\xi^3)\hat{z}(\xi)\yhwidehat{\varphi_2'z}(\xi) + (16\pi^4\xi^4)\yhwidehat{\varphi_2z}(\xi)\right]\mathrm{d}\xi 
\end{eqnarray*}
Defining $\lambda_3(\xi) := - \frac{i16\pi^3\xi^3}{(1+\xi^2)^3 }$ and $\lambda_4(\xi) :=  \frac{16\pi^4\xi^4}{(1+\xi^2)^3 }$ , we have that 
\[M_2z :=  \left(-\mathcal{F}^{-1}(\lambda_2)*\varphi_2''z\right)+\left(\mathcal{F}^{-1}(\lambda_3)*\varphi_2'z\right)+\left(\mathcal{F}^{-1}(\lambda_4)*\varphi_2z\right) \]
satisfies \fcref{c2}.

\section{aled}

Assuming that $f\in H^3(\R)$ is also in $L^1(\R)$, we can use the inversion formula and work in Fourier space :
\begin{equation}\label{int_finite}
K_n(f)(x) =  \sum_{i=1}^n\delta K(x,y_i)f(y_i) = \int_\R \hat{f}(\xi)\sum_{i=0}^n \delta K(x,y_i) e^{2i\pi\xi y_i}\mathrm{d}\xi
\end{equation}
On the other hand, the operator can be written as :
\begin{eqnarray*}
K(f)(x) &=& \int_\R K(x,y)f(y)\mathrm{d}y = \int_\R \hat{K}(x,\xi)\hat{f}(\xi)\mathrm{d}\xi \\
&=& \int_\R \hat{f}(\xi) \int_\R K(x,y)e^{-2i\pi\xi y}\mathrm{d}y \mathrm{d}\xi 
\end{eqnarray*}

There is a sign problem in the phase of the exponential, but if we can assume that the function $f$ is even, then we can also reverse the sign in the approximate kernel formula, and now the difference between the kernel and its approximation is easily bounded using the mean value theorem on $[-A,A]$, and by giving an explicit bound on the decay of the kernel away from zero. 

Now, taking the difference of the two expressions, we have :

\begin{eqnarray*}
\left| K(f)(x) - K_n(f)(x)\right| &\leq& \left|  \int_\R \hat{f}(\xi) \left(\int_{-A}^A K(x,y)e^{-2i\pi\xi y}\mathrm{d}y - \sum_{i=0}^n \delta K(x,y_i) e^{2i\pi\xi y_i} \right)\mathrm{d}\xi \right| \\ + \left|\int_\R \hat{f}(\xi) \int_{\R \setminus [-A,A]} K(x,y)e^{2i\pi\xi y}\mathrm{d}y\mathrm{d}\xi \right| \\
&\leq& (I) + (II)
\end{eqnarray*}

Let's first bound the difference on the compact domain $[-A,A]$ by applying the mean value theorem:

\begin{eqnarray*}
(I) &\leq&  \left|  \int_\R \hat{f}(\xi)  \sum_{i=0}^n \left(\int_{y_i}^{y_{i+1}} K(x,y)e^{2i\pi\xi y} - K(x,y_i) e^{2i\pi\xi y_i}\mathrm{d}y \right)\mathrm{d}\xi \right| \\
&\leq&   \int_\R \left| \hat{f}(\xi) \right|  \sum_{i=0}^n \left(\int_{y_i}^{y_{i+1}}  \delta \sup_{y\in[-A,A]} {\left|\partial_yK(x,y)e^{2i\pi\xi y} +2i\pi\xi K(x,y)e^{2i\pi\xi y} \right|} \mathrm{d}y \right)\mathrm{d}\xi \\
&\leq&  n  \delta^2  \int_\R \left| \hat{f}(\xi) \right|  \sup_{y\in[-A,A]}\Big( {\left|\partial_yK(x,y)\right|+ \left|2\pi\xi K(x,y)\right|}\Big) \mathrm{d}\xi \\
&\leq&  n  \delta^2 \Big(\| \hat{f}\|_{L^1}  \sup_{y\in[-A,A]}{\left|\partial_yK(x,y)\right|} +  \| \xi\hat{f}(\xi)\|_{L^1} \sup_{y\in[-A,A]}{2\pi|K(x,y)|} \Big) \\
&\leq&  C\frac{4A^2}{n}  \Big(\sup_{y\in[-A,A]}\left|\partial_yK(x,y)\right| + \sup_{y\in[-A,A]}\left|2\pi K(x,y)\right| \Big) \| f\|_{H^2} \\
\end{eqnarray*}

Where $C := \max\Big(\displaystyle \int_\R \frac{1}{(1+\xi^2)} \mathrm{d}\xi,\int_\R \frac{\xi^2}{(1+\xi^2)^2} \mathrm{d}\xi\Big)$ and we used the fact that, because $f\in H^3(\R)$, we have: \[\| \hat{f}(\xi)\|_{L^1} = \| \frac{1}{(1+\xi^2)^{1/2}}(1+\xi^2)^{1/2}\hat{f}(\xi)\|_{L^1} \leq \| \frac{1}{(1+\xi^2)^{1/2}}\|_{L^2} \|(1+\xi^2)^{1/2}\hat{f}(\xi)\|_{L^2} = \| \frac{1}{(1+\xi^2)^{1/2}}\|_{L^2} \| f\|_{H^1}  \] 
and  \[\| \xi\hat{f}(\xi)\|_{L^1} = \| \frac{\xi}{(1+\xi^2)}(1+\xi^2)\hat{f}(\xi)\|_{L^1} \leq \| \frac{\xi}{(1+\xi^2)}\|_{L^2} \|(1+\xi^2)\hat{f}(\xi)\|_{L^2} = \| \frac{\xi}{(1+\xi^2)}\|_{L^2} \| f\|_{H^2}  \] 

~\\
Then for the second part, by Fubini's theorem:
\begin{eqnarray*}
(II) &=&  \left|\int_{\R \setminus [-A,A]} K(x,y) \int_\R \hat{f}(\xi) e^{2i\pi\xi y}\mathrm{d}\xi \mathrm{d}y\right|  =\left| \int_{\R \setminus [-A,A]} K(x,y) f(y)\mathrm{d}y \right|  \\
&=& \left| \int_{\R \setminus [-A,A]} \lambda(x-y)\varphi(y) f(y)\mathrm{d}y \right| 
\end{eqnarray*}

We can compute that $\lambda(x)=\frac{\pi}{4}e^{-2\pi \lvert x\rvert}\left(2\pi^2x^2+3\pi \lvert x\rvert +\frac{3}{2}\right)$, and thus defining $P(x):=2\pi^2x^2+3\pi|x|+\tfrac32$ and recalling that  \(|x+y|\le|x|+|y|\) and \((x+y)^2\le2x^2+2y^2\), we can bound the lambda term:
\[
P(x+y)
\le2\pi^2(2x^2+2y^2)+3\pi(|x|+|y|)+\tfrac32
=2P(x)+2P(y).
\]
Hence
\[
\lambda(x+y)
=\tfrac\pi4e^{-2\pi|x+y|}P(x+y)
\le\tfrac\pi4e^{-2\pi|x+y|}(2P(x)+2P(y))
\]
\[
\leq 2e^{-2\pi|y|}\lambda(x)+2e^{-2\pi|x|}\lambda(y)
\]


Assuming we can show that $\varphi$ is bounded, we now have :
 We can also integrate by part here, since f is in $H^3$
\begin{eqnarray*}
(II) &\leq& \|\varphi\|_{L^\infty} \int_{\R \setminus [-A,A]} \Big( 2e^{-2\pi|y|}\lambda(x)+2e^{-2\pi|x|}\lambda(y)\Big)|f(y)|\mathrm{d}y \\
&\leq&  \|\varphi\|_{L^\infty} \left( \lambda(x) \int_{\R \setminus [-A,A]} 2e^{-2\pi|y|}|f(y)|\mathrm{d}y +   2e^{-2\pi|x|}\int_{\R \setminus [-A,A]} \lambda(y)|f(y)|\mathrm{d}y\right)\\
&\leq&  \|\varphi\|_{L^\infty} \left( \lambda(x) \left(\int_{\R \setminus [-A,A]} 2e^{-4\pi|y|}\mathrm{d}y\right)^\frac{1}{2} \|f\|_{L^2} +   2e^{-2\pi|x|}\int_{\R \setminus [-A,A]} \lambda(y)^2\mathrm{d}y\|f\|_{L^2}\right)\\
\end{eqnarray*}
\end{document}






























